\chapter{نتایج و بحث}

این فصل نتیجه پروژه را با زبان روشن بیان می‌کند: چه چیزهایی آماده است، چه چیزهایی هنوز نیاز به تکمیل دارد، و وضعیت کلی سامانه برای ارائه و توسعه بعدی چگونه است.

\section{ویژگی‌های پیاده‌سازی‌شده}
ویژگی‌های اصلی پروژه شامل کاتالوگ عمومی، صفحه جزئیات \lr{API}، مستندات، پلن‌ها، ثبت امتیاز، ثبت‌نام، ورود، نشست، خروج، پروفایل، چرخش \lr{API Key}، سازمان‌ها، اشتراک، \lr{checkout} دستی، مصرف، \lr{Caller}، \lr{Studio} و انتشار \lr{API} است. این ترکیب، پروژه را به فضای \lr{API Marketplace} نزدیک می‌کند \cite{rapidapi}.

این ویژگی‌ها دو دسته ارزش ایجاد می‌کنند. دسته نخست برای مصرف‌کننده \lr{API} است: کشف، مقایسه، مشاهده مستندات و انتخاب پلن. دسته دوم برای توسعه‌دهنده یا مالک \lr{API} است: مدیریت حساب، سازمان، کلید، انتشار، مشاهده مصرف و آزمایش جریان‌ها. ترکیب این دو دسته، پروژه را از یک فهرست ساده فراتر می‌برد و به شکل یک بازارچه عملیاتی نزدیک می‌کند.

\section{نمونه خروجی \lr{API}}
طرح \lr{OpenAPI} برای معرفی مسیرها و پاسخ‌ها استفاده می‌شود. نمونه زیر شکل ساده‌ای از پاسخ مسیر سلامت را نشان می‌دهد \cite{openapi}:
\begin{latin}
\begin{lstlisting}
GET /api/v1/system/health/
{
  "status": "ok",
  "timestamp": "2026-05-17T00:00:00Z"
}
\end{lstlisting}
\end{latin}

این پاسخ ساده است، اما از نظر ارزیابی اهمیت دارد؛ زیرا نشان می‌دهد مسیر نسخه‌دار \lr{API} فعال است و \lr{Client} می‌تواند وضعیت سرویس را بدون ورود بررسی کند. وجود زمان در پاسخ نیز برای عیب‌یابی و مانیتورینگ اولیه مفید است.

\section{نقاط قوت}
نقاط قوت پروژه عبارت‌اند از: مسیرهای نسخه‌دار، پاسخ خطای یکنواخت، آزمون‌های متعدد، پنهان‌کردن رازها، پشتیبانی از فارسی، قرارداد \lr{OpenAPI} و اجرای محلی با \lr{Docker Compose}.

از دید نگهداری، مهم‌ترین نقطه قوت پروژه جداشدن لایه‌هاست. مسیرها، \lr{Serializer}ها، لایه داده و \lr{Client} فرانت‌اند نقش‌های نسبتاً روشن دارند. از دید تجربه کاربر، پشتیبانی از راست‌به‌چپ و فارسی باعث می‌شود سامانه برای مخاطب ایرانی مناسب‌تر باشد. از دید توسعه، وجود آزمون‌ها و طرح \lr{OpenAPI} مسیر ادامه کار و اتصال \lr{Client}های دیگر را ساده‌تر می‌کند.

\section{چالش‌ها و بدهی فنی}
مهم‌ترین چالش، تصمیم نهایی درباره مرز داده و آماده‌سازی تولیدی است. تنظیمات توسعه‌ای باید از تنظیمات تولید جدا شود، رازها باید از محیط امن خوانده شوند، ورود اجتماعی به پیکربندی کامل نیاز دارد، و \lr{Caller} در نسخه تولیدی باید با کنترل خطا، محدودیت زمانی و ثبت رویداد دقیق به سرویس مقصد متصل شود.

این چالش‌ها مانع نمایش پروژه نیستند، اما برای بهره‌برداری واقعی باید حل شوند. دوگانگی داده می‌تواند در آینده باعث اختلاف رفتار بین پنل مدیریت و مسیرهای \lr{API} شود. تنظیمات توسعه‌ای نیز اگر بدون تغییر وارد محیط تولید شوند، خطر امنیتی ایجاد می‌کنند. همچنین پرداخت دستی و اجرای نمونه \lr{Caller} باید در نسخه تولیدی با سرویس‌های واقعی، کنترل خطا، محدودیت زمانی و ثبت رویداد دقیق جایگزین یا تکمیل شوند.

\section{بحث}
پیاده‌سازی فعلی برای نمایش و آزمون قابلیت‌های اصلی بازارچه \lr{API} مناسب است. با این حال، برای استفاده تولیدی باید مرز داده، محرمانگی تنظیمات، مدیریت سرویس واقعی \lr{MongoDB}، آزمون‌های فرانت‌اند و مسیر پرداخت واقعی دقیق‌تر شود.

سطح بلوغ پروژه را می‌توان «نمونه کامل قابل اجرا برای محیط توسعه» دانست. قابلیت‌های اصلی محصول وجود دارند و آزمون بک‌اند از آن‌ها پشتیبانی می‌کند، اما همه نیازهای یک سامانه تولیدی هنوز کامل نشده‌اند؛ از جمله مدیریت رازها، اتصال‌های بیرونی واقعی، مشاهده‌پذیری، پایداری زیر بار و چرخه انتشار رسمی.

\section{جمع‌بندی فصل}
نتیجه پروژه یک پورتال نسبتاً کامل برای کشف، انتشار و مدیریت \lr{API} است. بک‌اند از نظر آزمون‌های فعلی وضعیت خوبی دارد، اما آماده‌سازی تولیدی هنوز کار می‌خواهد. همین موضوع مبنای پیشنهادهای فصل پایانی است.
